%------------------------------------------------------------------------------%
\documentclass[fleqn,utf8,aspectratio=169,12pt]{beamer}

\usepackage{integracao_e_series}

\title{Resumo dos Testes de Convergência}

%------------------------------------------------------------------------------%
\begin{document}

\addtitle

%------------------------------------------------------------------------------%
\section{Resumo dos Testes}

%------------------------------------------------------------------------------%
\begin{frame}{Testes Básicos}
	\begin{itemize}[<+->]
		\item \emph{Teste da divergência ou do $n$-ésimo termo}\\
          a menos que $a_n\to 0$, a série diverge \\[5mm]
		\item \emph{Séries geométricas} \\[3mm]
          \(\sum_{n=1}^{\infty} ar^{n-1}\) converge se \(\abs{r}<1\); caso contrário, diverge \\[5mm]
		\item \emph{Série $p$} \\[3mm]
          \(\sum_{n=1}^{\infty} \frac{1}{n^p}\)
          converge se $p>1$; caso contrário, diverge\\[5mm]
    \item Série \emph{Telescópica}
	\end{itemize}
\end{frame}

%------------------------------------------------------------------------------%
\begin{frame}{Séries de Termos Não-Negativos}
\begin{itemize}[<+->]
  \item Teste da \emph{integral} \\[5mm]
  \item Teste da \emph{razão} ou o teste da \emph{raiz} \\[5mm]
  \item Teste da \emph{comparação}
\end{itemize}
\end{frame}

%------------------------------------------------------------------------------%
\begin{frame}{Séries Gerais -- Alguns Termos Negativos}
	\begin{itemize}[<+->]
    \item A série é \emph{alternada}? \\[3mm]
          Verifique as condições do teste da série alternada \\[5mm]
		\item \(\sum_{n=1}^{\infty} \abs{a_n}\) converge? \\[3mm]
		      Se convergir, \(\sum_{n=1}^{\infty} a_n\) também converge
	\end{itemize}
\end{frame}

%------------------------------------------------------------------------------%
\end{document}
%------------------------------------------------------------------------------%
